% Kurzfassung (Deutsch) und Abstract (Englisch)
\chapter*{Kurzfassung}
\addcontentsline{toc}{chapter}{Kurzfassung}

Reinforcement Learning wird zunehmend für autonome Cyberabwehr erprobt;
welche Rolle die Struktur des verteidigten Netzes dabei spielt, ist jedoch
wenig untersucht. Diese Arbeit misst den Einfluss der Netzwerktopologie
auf Verteidigungsagenten, die mit Proximal Policy Optimization in der
Simulationsumgebung CyberBattleSim trainiert werden, erweitert um einen
lernenden Verteidiger. Auf vier Topologien mit identischem
Inventar -- vollvermascht, Hub-and-Spoke, DMZ und mikrosegmentiert --
treten je sechs eingefrorene Angreifer an: fünf Spezialisten und ein über
alle Topologien gestuft trainierter Generalist. Die Datengrundlage bilden
120 Trainingsläufe über fünf Zufallsstartwerte sowie eine Evaluation mit
9000 Episoden gegen drei Vergleichsstufen. Ein Einfluss der Topologie auf
die Lerngeschwindigkeit ließ sich nicht nachweisen; auf das erreichte
Ergebnis ist er hingegen deutlich: Mit zunehmender Segmentierung sinkt der
Anteil der Episoden, in denen der Angreifer das wertvollste System
erreicht -- im mikrosegmentierten Netz bis auf null, während er im
vollvermaschten Netz nahezu ungebremst bleibt, weil Sperren dort
ausweichbar sind; erst Engpässe machen sie wirksam. Der auf der
verteidigten Topologie trainierte Angreifer erweist sich durchweg als
härtester Gegner. Methodisch zeigt die Arbeit, dass der Reward allein
Verteidigungsleistung nicht abbildet und das übernommene Reward-Design
Prävention unsichtbar ließ. Für die Praxis stützt das die Empfehlung
konsequenter Segmentierung: Sie begrenzt, was ein Angreifer erreichen
kann, und macht aktive Verteidigung erst wirksam.

\chapter*{Abstract}
\addcontentsline{toc}{chapter}{Abstract}

Reinforcement learning is increasingly explored for autonomous cyber
defense, yet the role played by the structure of the defended network has
received little attention. This thesis measures the influence of network
topology on defender agents trained with Proximal Policy Optimization in
the CyberBattleSim environment, extended with a learning defender. Four topologies with identical inventory -- fully meshed,
hub-and-spoke, DMZ, and micro-segmented -- are each defended against six
frozen attackers: five specialists and one generalist trained in stages
across all topologies. The data comprise 120 training runs over five
random seeds and an evaluation of 9{,}000 episodes against three baseline
tiers. No influence of topology on learning speed could be demonstrated;
its influence on the achieved outcome, however, is pronounced: With
increasing segmentation, the share of episodes in which the attacker
reaches the most valuable system declines -- down to zero in the
micro-segmented network, while in the fully meshed network the attacker
remains largely unchecked because blocks can simply be evaded; only choke
points make them effective. The attacker trained on the defended topology
consistently proves to be the hardest opponent. Methodologically, the
thesis shows that reward alone does not capture defensive performance and
that the inherited reward design left prevention invisible. For practice,
this supports a recommendation of rigorous segmentation: it limits what an
attacker can reach and is what makes active defense effective in the
first place.
